\chapter{Representación matricial y cambio de base}
\label{cap:representacion-matricial}

Una vez elegidas bases, cada transformación entre espacios de dimensión finita se representa mediante una matriz. Se estudia cómo dependen sus entradas de esas elecciones y por qué la semejanza expresa un cambio de representación \cite{lang1987,hoffmankunze1971,axler2024}.

\section{Matriz asociada a una transformación lineal}
Columnas como coordenadas de las imágenes de los vectores de la base.

\section{Composición y producto de matrices}
Compatibilidad entre composición de transformaciones y multiplicación matricial.

\section{Matriz de la transformación identidad}
Matriz identidad en una base y su comportamiento bajo cambios de coordenadas.

\section{Matriz de una transformación inversa}
Relación entre matrices inversas y transformaciones invertibles.

\section{Cambio de base en dominio y codominio}
Fórmulas de transformación de matrices cuando cambian las bases de entrada y salida.

\section{Matrices semejantes}
Definición, propiedades y ejemplos de matrices relacionadas por conjugación.

\section{Semejanza como cambio de representación}
Interpretación de matrices semejantes como representaciones de un mismo operador.
